\documentclass[12pt, a4paper]{article}

% ============================================================
% PACKAGES
% ============================================================
\usepackage[utf8]{inputenc}
\usepackage[T5]{fontenc}           % Tiếng Việt
\usepackage[vietnamese]{babel}     % Ngôn ngữ tiếng Việt
\usepackage{geometry}
\usepackage{titlesec}
\usepackage{booktabs}
\usepackage{array}
\usepackage{tabularx}
\usepackage{xcolor}
\usepackage{hyperref}
\usepackage{enumitem}
\usepackage{fancyhdr}
\usepackage{graphicx}
\usepackage{amsmath}
\usepackage{cite}

% ============================================================
% PAGE SETUP
% ============================================================
\geometry{
    top=2.5cm,
    bottom=2.5cm,
    left=3cm,
    right=2cm
}

% ============================================================
% COLORS
% ============================================================
\definecolor{primaryblue}{RGB}{0, 70, 127}
\definecolor{accentred}{RGB}{192, 0, 0}
\definecolor{lightgray}{RGB}{245, 245, 245}
\definecolor{darkgray}{RGB}{80, 80, 80}

% ============================================================
% HYPERLINK SETUP
% ============================================================
\hypersetup{
    colorlinks=true,
    linkcolor=primaryblue,
    urlcolor=primaryblue,
    citecolor=accentred
}

% ============================================================
% HEADER & FOOTER
% ============================================================
\pagestyle{fancy}
\fancyhf{}
\fancyhead[L]{\small\textcolor{darkgray}{Nghiên cứu NLP -- Xử lý Ngôn ngữ Tự nhiên}}
\fancyhead[R]{\small\textcolor{darkgray}{\thepage}}
\fancyfoot[C]{\small\textcolor{darkgray}{\textit{Báo cáo môn học -- \the\year}}}
\renewcommand{\headrulewidth}{0.4pt}

% ============================================================
% SECTION FORMATTING
% ============================================================
\titleformat{\section}
    {\large\bfseries\color{primaryblue}}
    {\thesection.}{0.5em}{}[\titlerule]

\titleformat{\subsection}
    {\normalsize\bfseries\color{darkgray}}
    {\thesubsection.}{0.5em}{}

% ============================================================
% DOCUMENT
% ============================================================
\begin{document}

% ------------------------------------------------------------
% TITLE PAGE
% ------------------------------------------------------------
\begin{titlepage}
    \centering
    \vspace*{1cm}

    {\Large\bfseries TRƯỜNG ĐẠI HỌC [TÊN TRƯỜNG]}\\[0.3cm]
    {\large KHOA [TÊN KHOA]}\\[2cm]

    \rule{\linewidth}{1.5pt}\\[0.4cm]
    {\LARGE\bfseries\color{primaryblue}
        BÁO CÁO NGHIÊN CỨU\\[0.3cm]
        XỬ LÝ NGÔN NGỮ TỰ NHIÊN\\
        (NATURAL LANGUAGE PROCESSING)
    }\\[0.4cm]
    \rule{\linewidth}{1.5pt}\\[2cm]

    \begin{tabular}{ll}
        \textbf{Môn học:}       & [Tên môn học] \\[0.3cm]
        \textbf{Giảng viên:}    & [Tên giảng viên] \\[0.3cm]
        \textbf{Nhóm:}          & [Số nhóm] \\[0.3cm]
        \textbf{Thành viên:}    & [Họ tên thành viên 1] \\
                                & [Họ tên thành viên 2] \\
                                & [Họ tên thành viên 3] \\[0.3cm]
        \textbf{Năm học:}       & 2024 -- 2025 \\
    \end{tabular}

    \vfill
    {\large\the\year}
\end{titlepage}

% ------------------------------------------------------------
% TABLE OF CONTENTS
% ------------------------------------------------------------
\tableofcontents
\newpage

% ============================================================
% SECTION 1: ĐỀ TÀI NGHIÊN CỨU
% ============================================================
\section{Đề Tài Nghiên Cứu}

\subsection{Tên đề tài}

\begin{center}
    \colorbox{lightgray}{%
        \parbox{0.9\textwidth}{%
            \centering\vspace{0.3cm}
            {\large\bfseries Xử lý Ngôn ngữ Tự nhiên (Natural Language Processing -- NLP)}\\
            \vspace{0.1cm}
            \textit{Ứng dụng mô hình Transformer vào phân loại văn bản tiếng Việt}
            \vspace{0.3cm}
        }%
    }
\end{center}

\subsection{Lý do chọn đề tài}

Xử lý ngôn ngữ tự nhiên (NLP) là một lĩnh vực quan trọng trong trí tuệ nhân tạo, 
đặc biệt với sự bùng nổ của các mô hình ngôn ngữ lớn như BERT, GPT và các biến thể. 
Đối với tiếng Việt -- một ngôn ngữ có đặc thù về thanh điệu và cấu trúc -- 
việc nghiên cứu và ứng dụng NLP còn nhiều tiềm năng phát triển.

% ============================================================
% SECTION 2: CÂU HỎI NGHIÊN CỨU
% ============================================================
\section{Câu Hỏi Nghiên Cứu (Research Question)}

\subsection{Research Question (RQ)}

\begin{center}
    \fbox{%
        \parbox{0.85\textwidth}{%
            \vspace{0.3cm}
            \centering
            \textit{``What transformer-based NLP models have been applied to Vietnamese 
            text classification in the last 3 years?''}
            \vspace{0.3cm}
        }%
    }
\end{center}

\vspace{0.5cm}
\noindent\textbf{Dịch nghĩa:} Những mô hình NLP dựa trên kiến trúc Transformer nào 
đã được ứng dụng vào bài toán phân loại văn bản tiếng Việt trong 3 năm gần đây?

\subsection{Phân tích RQ theo khung PICO}

\begin{table}[h!]
\centering
\begin{tabularx}{0.9\textwidth}{>{\bfseries}l X}
\toprule
Thành phần & Nội dung \\
\midrule
\textbf{P}opulation    & Văn bản tiếng Việt (tin tức, mạng xã hội, tài liệu) \\
\textbf{I}ntervention  & Mô hình Transformer (BERT, PhoBERT, XLM-R, \ldots) \\
\textbf{C}omparison    & Các phương pháp truyền thống (SVM, LSTM, CNN, \ldots) \\
\textbf{O}utcome       & Độ chính xác phân loại văn bản (Accuracy, F1-score) \\
\bottomrule
\end{tabularx}
\caption{Phân tích câu hỏi nghiên cứu theo khung PICO}
\end{table}

% ============================================================
% SECTION 3: TỪ KHÓA
% ============================================================
\section{Từ Khóa (Keywords)}

\begin{table}[h!]
\centering
\begin{tabularx}{0.85\textwidth}{c >{\bfseries}l X}
\toprule
\textbf{STT} & \textbf{Từ khóa} & \textbf{Giải thích} \\
\midrule
1 & Natural Language Processing (NLP) 
  & Lĩnh vực AI liên quan đến xử lý và hiểu ngôn ngữ của con người \\
\addlinespace
2 & Transformer / BERT 
  & Kiến trúc mạng nơ-ron tiên tiến được sử dụng rộng rãi trong NLP \\
\addlinespace
3 & Vietnamese Text Classification 
  & Bài toán phân loại văn bản áp dụng cho ngôn ngữ tiếng Việt \\
\bottomrule
\end{tabularx}
\caption{Ba từ khóa chính của nghiên cứu}
\end{table}

\noindent\textbf{Chuỗi tìm kiếm gợi ý:}
\begin{quote}
    \texttt{("NLP" OR "Natural Language Processing") AND ("Transformer" OR "BERT") 
    AND ("Vietnamese" OR "text classification") AND (2022:2025)}
\end{quote}

% ============================================================
% SECTION 4: TIÊU CHÍ CHỌN / LOẠI BÀI
% ============================================================
\section{Tiêu Chí Lựa Chọn Nghiên Cứu}

\subsection{Tiêu chí chọn bài (Inclusion Criteria)}

\begin{enumerate}[leftmargin=*, label=\textbf{I\arabic*.}]
    \item Bài báo xuất bản từ năm \textbf{2022 đến nay} (3 năm gần đây).
    \item Nghiên cứu liên quan đến \textbf{NLP cho tiếng Việt} hoặc ngôn ngữ có đặc thù tương tự.
    \item Có sử dụng \textbf{mô hình học sâu} (deep learning) hoặc kiến trúc Transformer.
    \item Được đăng trên \textbf{tạp chí/hội nghị khoa học} có peer-review (ACL, EMNLP, IEEE, Scopus\ldots).
    \item Có \textbf{kết quả thực nghiệm} với dataset cụ thể và metrics đánh giá rõ ràng.
\end{enumerate}

\subsection{Tiêu chí loại bài (Exclusion Criteria)}

\begin{enumerate}[leftmargin=*, label=\textbf{E\arabic*.}]
    \item Bài viết \textbf{không phải tiếng Anh hoặc tiếng Việt}.
    \item Nghiên cứu chỉ dùng \textbf{phương pháp truyền thống} (rule-based, không có ML/DL).
    \item Bài \textbf{trùng lặp} hoặc chỉ là abstract (không có full-text).
    \item \textbf{Không có kết quả thực nghiệm} hoặc không so sánh với baseline.
    \item Bài thuộc \textbf{lĩnh vực khác} (xử lý ảnh, âm thanh đơn thuần không kết hợp NLP).
\end{enumerate}

\subsection{Tổng hợp tiêu chí (PRISMA-style)}

\begin{table}[h!]
\centering
\begin{tabularx}{0.9\textwidth}{l X X}
\toprule
\textbf{Tiêu chí} & \textbf{Chọn (Inclusion)} & \textbf{Loại (Exclusion)} \\
\midrule
Thời gian    & 2022 -- nay              & Trước 2022 \\
Ngôn ngữ     & Tiếng Anh / Việt         & Ngôn ngữ khác \\
Phương pháp  & Deep Learning, Transformer & Rule-based thuần túy \\
Nguồn        & Peer-reviewed journal/conf & Blog, báo cáo không rõ nguồn \\
Nội dung     & Có thực nghiệm, metrics   & Chỉ có abstract, trùng lặp \\
\bottomrule
\end{tabularx}
\caption{Bảng tổng hợp tiêu chí lựa chọn nghiên cứu}
\end{table}

% ============================================================
% SECTION 5: KẾ HOẠCH THỰC HIỆN (tùy chọn)
% ============================================================
\section{Kế Hoạch Thực Hiện}

\begin{enumerate}
    \item \textbf{Tuần 1--2:} Tìm kiếm tài liệu trên Google Scholar, Semantic Scholar, ACL Anthology.
    \item \textbf{Tuần 3--4:} Sàng lọc theo tiêu chí Inclusion/Exclusion.
    \item \textbf{Tuần 5--6:} Đọc và trích xuất thông tin từ các bài được chọn.
    \item \textbf{Tuần 7--8:} Tổng hợp, phân tích và viết báo cáo.
\end{enumerate}

% ============================================================
% TÀI LIỆU THAM KHẢO (placeholder)
% ============================================================
\section*{Tài Liệu Tham Khảo}
\addcontentsline{toc}{section}{Tài Liệu Tham Khảo}

\begin{thebibliography}{9}

\bibitem{devlin2019bert}
Devlin, J., Chang, M. W., Lee, K., \& Toutanova, K. (2019).
\textit{BERT: Pre-training of deep bidirectional transformers for language understanding}.
NAACL-HLT 2019.

\bibitem{phobert2020}
Nguyen, D. Q., \& Nguyen, A. T. (2020).
\textit{PhoBERT: Pre-trained language models for Vietnamese}.
Findings of EMNLP 2020.

\bibitem{vaswani2017attention}
Vaswani, A., et al. (2017).
\textit{Attention is all you need}.
NeurIPS 2017.

\end{thebibliography}

\end{document}